\section{SAE Analysis}
\label{sec:sae_analysis}

This section separates diagnostics already available from the completed study from proposed analyses of the intervention mechanism.
The existing evidence establishes length associations and shorter generated outputs; it does not establish that the selected features encode redundant reasoning.
The analyses below are intended to test that interpretation and to identify which aspects of the intervention merit downstream evaluation.

\subsection{Existing Evidence: Answer-Marker Associations}
\label{sec:answer_marker_evidence}

The held-out token-context report for the primary SAE shows that the highest-activation examples of all eight selected short-associated features occur at \texttt{Answer} tokens.
For example, feature F7339 has high-activation contexts such as a completed calculation followed by \texttt{Answer:}; analogous contexts occur in both short and long traces.
These examples indicate an answer-format association, but do not measure how much of the feature's total activation is explained by the marker.
They motivate distinguishing a feature that tracks useful concise reasoning from one that tracks answer formatting or the transition into an answer segment.

A simple aggregation effect can also produce a short--long association.
If a marker appears once per trace and contributes approximately constant activation $A_j$, a feature's whole-trace mean can behave as
\begin{equation}
  \bar z_j(\tau)=\frac{1}{L_\tau}\sum_{t=1}^{L_\tau} z_{t,j}
  \approx \frac{A_j}{L_\tau}.
\label{eq:length_denominator}
\end{equation}
Thus, token-normalized activation can correlate with shortness even when activation at the marker is unchanged.
This is a possible explanation, not a fitted result.
The existing first-64-token direction check does not by itself resolve it, because some traces may already contain an answer marker within that window.

\subsection{Separating Length, Position, and Lexical Cues}
\label{sec:sae_falsification}

We propose decomposing feature activation into reasoning-body, answer-marker, and answer-suffix regions, reporting both region-specific means and contributions to total activation.
The analysis would compare short and long traces within the same problem and, where supported, at matched token identities and absolute positions.
A before/after comparison excluding answer regions would test whether the observed association survives removal of its most evident correlate.
Regressing whole-trace means against inverse length provides a complementary diagnostic of Equation~\ref{eq:length_denominator}; it does not replace the within-problem analysis.

A controlled text battery would preserve mathematical steps while changing wording or answer headings, preserve formatting while changing calculation correctness, and inject associated tokens into non-reasoning text~\citep{ma2026reasoning}.
Step and answer checks are needed to verify the intended manipulations.
Because deletion and rewriting shift token positions and subsequent states, analysis would align corresponding spans and include position-matched controls rather than compare only whole-trace averages.
Passing these tests would strengthen a semantic interpretation; failing them could still leave a useful formatting-control direction.

Early-prefix prediction would use only information available after 32 or 64 generated tokens to predict remaining length or subsequent entry into the answer segment.
Evaluation would group by problem and report the fraction of traces ending before each landmark.
Relative positions computed using eventual trace length and retrospective end-of-trace masks would be restricted to offline diagnostics, avoiding future information in an online predictor.

\subsection{Target Engagement and Direction Specificity}
\label{sec:sae_engagement}

The selected intervention adds a constant composite decoder direction scaled to the current hidden-state norm.
To determine what it actually changes, we propose encoding the same states before and after the implemented perturbation:
\begin{equation}
  z_t=\operatorname{Enc}(h_t),\qquad
  z'_t=\operatorname{Enc}(h_t+\delta_t).
\label{eq:feature_readback}
\end{equation}
Readback must reproduce the deployed layer, normalization, feature aggregation, and numerical precision.
Comparing independently generated texts would confound the perturbation with changed token histories.

Measurements would include the signed change in selected features, the magnitude of changes in unselected features, turnover of the TopK active set, and actual perturbation-to-state norm.
The true next-token distribution after the remaining model layers would provide KL divergence and probabilities of answer markers and end-of-sequence tokens.
These quantities are especially relevant at $\rho=0.30$: matching hidden-state perturbation norms does not ensure comparable distribution shifts or feature specificity.
A fixed decoder-direction injection and activation-gated feature suppression would remain separately named interventions.

\subsection{Feature Redundancy, Timing, and Stability}
\label{sec:sae_structure}

\paragraph{Composition of the eight-feature direction.}
We propose inspecting pairwise decoder cosine similarity, activation correlations, and the singular-value spectrum of the selected directions.
Shared answer-marker contexts make it important to determine whether eight features provide several independent effects or largely overlapping directions.
Single-feature, top-1/2/4/8, and leave-one-out interventions would hold the total perturbation norm fixed.
Directions would be ranked on development data before downstream evaluation.
Larger sets would require additional features passing the same discovery rule, rather than relaxing selection thresholds to fill a requested set size.

\paragraph{When shortening occurs.}
Matched continuations from a shared natural prefix would compare enhancement, zero intervention, sign reversal, and multiple random controls with independent caches.
Early-only, delayed, and answer-marker-triggered windows would localize the effect using only information available during generation.
Outputs would be decomposed into answer-preceding steps, repeated explanations, and answer-following text.
A reduction concentrated after the answer would support a different interpretation from preserved calculations with fewer repeated derivations.
Response correctness and usable-answer coverage would accompany every length comparison.

\paragraph{Dictionary quality and reproducibility.}
The six existing SAEs support a first comparison across layers and sparsities without training a new grid.
Their reconstruction and activation-coverage metrics are reported in Appendix~\ref{app:sae_quality}.
Additional seeds for the primary setting would test whether the feature structure and intervention effects recur under dictionary retraining.
Matching would use decoder geometry and activation profiles on independent questions, allowing feature splitting; feature IDs and the number of significant candidates are not measures of semantic stability.

\subsection{Connecting Feature Changes to Student Learning}
\label{sec:sae_student_link}

The decisive endpoint remains student performance after training on the generated data.
We propose comparing SAE, dense or ASC, and prompt-based traces at similar length distributions under the training controls in Section~\ref{sec:baseline_design}.
Method-blinded annotation would distinguish calculations, explanations, verification, repeated content, and answer formatting on paired problems.
Student negative log-likelihood over these segments could provide a low-cost diagnostic of learnability, but lower likelihood loss would not establish a subsequent generalization gain.
The multi-answer extension in Appendix~\ref{app:scope} would additionally separate multiple distinct texts from repeated exposure to the same text, subject to its unequal-token limitation.
These analyses would connect a measured change in teacher behavior to the content and learning effects of the resulting supervision, without presupposing an accuracy benefit.
